The index has three lifecycle stages: training produces reusable model
state, adding encodes and physically arranges database rows, and searching
uses the completed arrays. This separation matters experimentally. A
trained-state cache avoids learning the model again, but still requires
database encoding and list construction. A cold train-plus-add run and a
search over a reused state measure different operations.

\begin{figure}[tbp]
\centering
% Audited lifecycle; no proposed storage or distributed subsystem.
\begin{tikzpicture}[
 font=\small\linespread{1}\selectfont,
 box/.style={draw=ThesisNavy!65,rounded corners=2pt,align=center,
   text width=3.65cm,minimum height=1.9cm,inner sep=5pt},
 gpu/.style={box,fill=ThesisNavy!7},
 arr/.style={-{Latex[length=2mm]},draw=ThesisNavy!80,thick}]
\node[box] (host) at (0,0) {Host input vectors\\Training samples\\Database batches};
\node[box] (model) at (4.8,0) {Host model\\QR rotation $\Pi$\\Primary codebook};
\node[box] (cache) at (9.6,0) {Host model cache\\$\Pi$ and codebooks\\IVF centroids};
\node[gpu] (train) at (0,-3.6) {GPU training\\Rotate samples\\Fit IVF centroids\\Fit residual model};
\node[gpu] (add) at (4.8,-3.6) {GPU bulk build\\1 or 2 passes\\Assign / encode\\Reorder / pack\\Floats $\leq768$ MiB};
\node[gpu] (index) at (9.6,-3.6) {GPU-resident index\\Offsets and IDs\\Primary / residual codes\\Corrections};
\node[box] (query) at (0,-7.2) {Host query batch\\Pinned staging\\Transfer to GPU};
\node[gpu] (search) at (4.8,-7.2) {GPU query pipeline\\Rotate; probe; scan\\Refine; final top-$k$\\Workspace / graph};
\node[box] (result) at (9.6,-7.2) {Host query results\\IDs and scores\\Transfer from GPU};
\draw[arr] (host) -- node[left,font=\footnotesize] {samples} (train);
\draw[arr] (model.south west) -- (train.north east);
\draw[arr] (train) -- node[above,font=\footnotesize] {model} (add);
\draw[arr] (host.west) -- ++(-.5,0) |- (4.8,-5.15)
 node[pos=.8,below,font=\footnotesize] {database batches} -- (add.south);
\draw[arr] (model) -- (add);
\draw[arr] (add) -- (index);
\draw[<->,dashed,draw=ThesisNavy!80,thick]
 (train.north east) -- node[pos=.8,above,sloped,font=\footnotesize,fill=white,inner sep=1pt]
 {save / load model} (cache.south west);
\draw[arr] (query) -- (search);
\draw[arr] (search) -- (result);
\draw[arr] (index.south) -- ++(0,-.6) -| (search.north);
\end{tikzpicture}

\caption[Build and query architecture]{Implemented build and query boundaries. Host work includes
sampling, orthogonal-matrix construction and analytical primary-codebook
generation; the evaluated configuration fits IVF and residual models on
the GPU. The cache stores trained state, whereas the completed database
representation resides on the device. Dashed arrows denote trained-state
reuse. The staging label bounds only raw/rotated float buffers at source
defaults; other pipeline scratch is additional.}
\label{fig:4-architecture}
\end{figure}

The host owns the input vectors and supplies training samples. It constructs
the square transform and analytical primary codebook. GPU matrix products
apply the transform; the device training path fits coarse IVF centroids
and scalar residual codebooks. Adding then transfers bounded batches,
encodes both levels, determines IVF membership, and writes a common
list-ordered representation. Host orchestration and device computation
therefore coexist within a GPU-resident search system.

The model state consists of the transform, primary codebook, residual
codebooks and IVF centroids. The database state adds list offsets, original
row identifiers, primary codes, residual codes and correction factors.
Search workspace has a shorter logical lifetime: rotated queries, centroid
scores, probe offsets, LUTs and survivor buffers are overwritten on the
next batch, although their allocations are retained. A captured CUDA graph
belongs to this execution state and references the existing arrays.

The audited API builds a fixed snapshot with one successful \texttt{add}
call; it rejects a subsequent add. Physical positions remain fixed while
searching that snapshot. Changes are handled by building a replacement snapshot. This gives a
simple bulk-rebuild lifecycle; mutable insert/delete operations and
transactional publication would require additional index management.

Search accepts host queries and returns host IDs and computed scores.
The query-only constant in those scores and its implications for range
thresholds are derived in \cref{eq:5-refined-constant}. The measured
host-to-host call boundary is defined in \cref{sec:3-3}.


This design streams construction inputs from host memory while retaining
the compressed database on the GPU. It supplies neither a disk-backed
search path nor a claim that an arbitrarily large compressed index fits
one device.
